\documentclass[notitlepage,onecolumn,showpacs,11pt]{revtex4-1}
\usepackage{amssymb}
\usepackage{graphicx}
\usepackage{amsmath}
\usepackage{float}
\usepackage{epstopdf}
\usepackage{times}
\usepackage{color}
\usepackage{subfigure}
\usepackage{setspace}
\usepackage{bm}% bold math

\newcommand{\bsym}[1]{\boldsymbol{#1}}
\newcommand{\mbf}[1]{\mathbf{#1}}
\newcommand{\grad}{\boldsymbol{\nabla}}
\newcommand{\Tr}{\text{Tr}}

\begin{document}

\title{Bond-disordered non-collinear magnets: Stability of the goldstone modes}

\maketitle

Much like temperature, in some systems quenched disorder prevents spontaneous 
symmetry breaking below a certain lower critical dimension \cite{imry-ma-75}. 
For continuous symmetries, break down of a long ranged order is signalled by 
the infrared catastrophe of the Goldstone mode fluctuations of the ordered state. 
This argument, attributed to Mermin and Wagner \cite{mermin-wagner-66}, has a 
straightforward parallel in magnetically ordered systems coupled with random 
fields \cite{aharony-78}.  In the following we shall show that the random 
exchange couplings also become crucial to the fate of a certain class of magnetic 
order, the non-collinear antiferromagnets.

\section{The goldstone mode theory of collinear and non-collinear magnets}

\subsection{Goldstone modes in collinear magnets}

\subsubsection{The model}

Massless goldstone modes are generated when the continuous symmetry of a hamiltonian
is spontaneously broken. The classic example is an exchange coupled ferromagnet with 
a global spin rotation invariance. The continuum approximation for the hamiltonian 
for such a system is given by, 
\begin{equation}
    \begin{aligned}
        H = \frac{\rho}{2}\int d^dx 
        \left(\nabla\mbf{S}(x)\right)^2
    \end{aligned}
\end{equation}
where $\rho$ is the spin-stiffness parameter. At zero temperature the ground state 
breaks the full $O(3)$ rotational symmetry of the model and assumes a configuration 
with all spins parallel to each other in a collinear pattern. The ground state 
therefore has a reduced $O(2)$ symmetry related to the rotation around the magnetization 
axis. In a generic model with a symmetry breaking from $O(N)$ to $O(N-1)$, the elementary 
excitations are the goldstone modes linked with the generators of the quotient group 
$O(N)/O(N-1)$, also called, the broken generators. Let us consider a symmetry broken 
ground state for an $O(N)$ ferromagnet, $\mbf{S}_0 = \left(0 ... 0, 1\right)^T $. It is
immediately apparent that there is a degenerate continuum of states related to the 
residual $O(N-1)$ rotational symmetry spanning the $N-1$ dimensional subspace orthogonal 
to the magnetization axis.

Any symmetry operation that goes beyond this subspace would generate massless 
goldstone modes. We can make a particular gauge choice,

\begin{equation}
    \begin{aligned}
        U =  \begin{pmatrix}
            \sqrt{I_{N-1}-\bsym{\pi}^T\bsym{\pi}} 
            & \bsym{\pi}\\
            -\bsym{\pi}^T
            & \sqrt{1-\bsym{\pi}\bsym{\pi}^T} 
        \end{pmatrix}
    \end{aligned}
\end{equation}
where, $U\in O(N)/O(N-1)$ and find, $\mbf{S}=U\mbf{S}_0=(\bsym{\pi},
\sqrt{1-\bsym{\pi}\bsym{\pi}^T})^T$. By pluggin this back to the hamiltonian
we immediately arrive at the familiar $O(N)$ NL$\sigma$-model hamiltonian.
\begin{equation}
    \begin{aligned}
        H &= \frac{\rho}{2}\int d^dx\left(\grad\pi_i\cdot\grad\pi_i
        +\frac{(\pi_i\grad\pi_i)^2}{1-\pi_i\pi_i}\right)\\
        &= \int d^dx g_{ij}(\bsym{\pi})\grad\pi_i\cdot\grad\pi_j
    \end{aligned}
\end{equation}
For classical spins, this model is renormalizable around the lower critical
dimension $d=2+\epsilon$ with dimensional regularization. The metric $\beta g_{ij}$ 
effectively plays the role of an inverse temperature and the perturbative
expansion can be organized as a low temperature expansion. For quantum spins
the action takes up the form,
\begin{equation}
    \begin{aligned}
        S = \int d\tau\int d^dx g_{ij}(\bsym{\pi})
        \partial_\mu\pi_i\partial^\mu\pi_j
    \end{aligned}
\end{equation}
and a similar renormalization scheme follows at $d=1+\epsilon$. Although, in this
case the metric is understood to be the stiffness parameter for the spins. 

The $O(N)$ NL$\sigma$-model is however special because although the expansion of the 
metric coefficient gives rise to infinite number of vertices, there is only one coupling 
that flows (here, $\rho$). Yet, NL$\sigma$-model is renormalizable at $D=d+z=2$ and upto
one-loop order one can find the renormalized value of the stiffness parameter or 
temperature by considering the perturbative correction to the propagator alone. Given
the model is massless, in a wilsonian treatment the only non-trivial element turns out 
to be the field strength renormalization. A standard technique to resolve this
problem is to consider a trivial renormalization of an external field coupled to the
spin or alternatively ensuring that both the massless $\bsym{\pi}$ goldstone modes and
the massive $\sigma=\sqrt{1-\bsym{\pi}\bsym{\pi}^T}$ mode have the same scaling. Both
of these schemes are manifestation of the fact that one has to consider the fully symmetric 
model to find the scaling parameters.

\subsubsection{Renormalization at D = 2}

\subsection{NL$\sigma$-model in a symmetric coset space}

\subsubsection{The model}

Although, the goldstone mode effective action of the $O(N)$ symmetric ferromagnet
followed quite naturally by considering a particular gauge choice for the broken generators
in the preceeding section, in the following we shall consider a more generic approach
to arrive at a NL$\sigma$-model effective action of a spontaneous symmetry breaking 
$G/H$. This will be useful for the particular case of our interest, the non-collinear 
antiferromagnets. In the context of the following discussion, e.g. $G/H=O(N)/O(N-1)$ shall 
describe what is known as a symmetric coset space but the extension to the non-collinear case
will be more complicated than that.

\subsubsection{The renormalization}

\subsection{NL$\sigma$-model for non-collinear antiferromagnets}

\subsubsection{The model}

As we saw in the preceeding section, the $O(N)/O(N-1)$ spontaneous symmetry breaking of 
simple ferromagnets can be understood in the context of symmetric coset space and a 
single coupling parameter, identified as either the the spin-stiffness or the termperature
depending on whether one considers quantum spins or finite temperature classical spins. The
situation is slightly different with the non-collinear magnets, e.g. the triangular lattice 
heisenberg antiferromagneti where the original symmetry of the spins are fully broken 
with $O(3)/Z_2\simeq SO(3)$ but there is still a residual $C_{3v}$ symmetry of the lattice
left in the problem. In a continuum description, this leads to a symmetry breakinig pattern
of $O(3)\times O(2)/O(2)$, where the $O(2)$ is not a physical symmetry of the goldstone
bosons but the extension of the lattice symmetry \cite{azaria-93}. As we show below, it is
possible to arrive at this conclusion from a microscopic analysis \cite{dombre-88}. 

We consider the antiferromagnetic heisenberg hamilonian in a triangular lattice and consider 
an expansion around the $120^\circ$ ordered state. The goldstone modes shall be excited by
the generators of the broken symmetry group, $SO(3)$. Our microscopic to continuum modelling 
would involve considering plaquettes of elementary upward triangles, finding the low energy
contributions of a departure from the $120^\circ$ order and ultimately coarse-graining within 
each elementary triangle. We start by denoting the three sublattice orientaions of the 
$120^\circ$ order, $\mbf{n}_A,\mbf{n}_B$ and $\mbf{n}_C = -(\mbf{n}_A+\mbf{n}_B)$ with a
spin-excitation given by $\mbf{S}_\alpha=U\mbf{n}_\alpha$, where $U\in SO(3)$. We consider the
hamiltonian within an elementary triangular plaquette.
\begin{equation}
    \begin{aligned}
        H = J \left(\mbf{S}_A(\mbf{x}_A)\cdot\mbf{S}_B(\mbf{x}_B)
        +\mbf{S}_B(\mbf{x}_B)\cdot\mbf{S}_C(\mbf{x}_C)
        +\mbf{S}_C(\mbf{x}_C)\cdot\mbf{S}_A(\mbf{x}_A)
        \right) 
    \end{aligned}
\end{equation}
By performing a gradient expansion coupled with a coarse-graining of the spatial co-ordinate to 
capture the low energy pieces, we arrive at the following expression.
\begin{equation}
    \begin{aligned}
        H &= JS^2a^2 \left(\mbf{n}_A U^{-1}\cdot\grad^2U\mbf{n}_B
        -\mbf{n}_B U^{-1}\cdot\grad^2U(\mbf{n}_A+\mbf{n}_B)
        -\mbf{n}_A U^{-1}\cdot\grad^2U(\mbf{n}_A+\mbf{n}_B)
        \right)\\
        &= JS^2a^2 \left(\mbf{n}_A U^{-1}\cdot\grad^2U\mbf{n}_B
        -(\mbf{n}_A+\mbf{n}_B)U^{-1}\cdot\grad^2U(\mbf{n}_A+\mbf{n}_B)
        \right) 
    \end{aligned}
\end{equation}
The single gradient term has been dropped from the expansion as it can be easily shown to have
no contribution at all from the analysis that follows. Now, a transformation to the orthonormal basis, 
$\mbf{e}_1=\sqrt{\frac{3}{2}}\left(\mbf{n}_A+\mbf{n}_B\right)$ and 
$\mbf{e}_2=\sqrt{\frac{1}{2}}\left(\mbf{n}_A-\mbf{n}_B\right)$ re-arranges the terms of the 
hamiltonian in a convenient form.
\begin{equation}
    \begin{aligned}
        H &= -\frac{JS^2a^2}{2} \left(\mbf{e}_1 U^{-1}\cdot\grad^2U\mbf{e}_1
        +\mbf{e}_2 U^{-1}\cdot\grad^2U\mbf{e}_2
        \right)
    \end{aligned}
\end{equation}
An extra $O(2)$ symmetry within the $(\mbf{e}_1,\mbf{e}_2)$ vector space is immediately apparent.
Defining $S_0=(\mbf{e}_1,\mbf{e}_2)$ and performing a gauge transformation $U = US_0$ we can re-write,
\begin{equation}
    \begin{aligned}
        H &= -\frac{JS^2a^2}{2} \Tr\left(S_0^{-1} U^{-1}\grad^2U S_0\right)\\
        &= -\frac{JS^2a^2}{2} \Tr\left(U^{-1}\grad^2U\right)
    \end{aligned}
\end{equation}
However, it has been shown that for a complete renormalization of the problem one has to consider a
third unit vector $\mbf{e}_3=\mbf{e}_1\times\mbf{e}_{2}$ \cite{azaria-93}. This introduces a second
coupling to the problem with the complete hamiltonian now given by,
\begin{equation}
    \begin{aligned}
        H &= p_1\left(\mbf{e}_1 U^{-1}\cdot\grad^2U\mbf{e}_1
        +\mbf{e}_2 U^{-1}\cdot\grad^2U\mbf{e}_2
        \right)
        +p_2\left(\mbf{e}_3 U^{-1}\cdot\grad^2U\mbf{e}_3\right)\\
        &=\Tr\left(P(U^{-1}\grad U)^2\right)
    \end{aligned}
\end{equation}
with gauge transformation $U = U (\mbf{e}_1,\mbf{e}_2,\mbf{e}_3)$ generating the full $SO(3)$
goldstone modes and the diagonal coupling matrix $P=\text{diag}(p_1,p_1,p_2)$ ensuring the 
$O(3)\times O(2)/O(2)$ nature of the exciatations. Given $\pi_1(SO(3))=Z_2$, the order parameter 
$U$ can have $Z_2$ point defects which are gapped out at zero temperature. With this knowledge we 
can immediately write down the quantum action of the full goldstone mode problem 
\cite{dombre-88,azaria-92}.
\begin{equation}
    \begin{aligned}
        S &=\int d\tau\int d^dx 
        \left[
        \Tr\left(P_\tau(U^{-1}\partial_\tau U)^2\right)
        +
        \Tr\left(P_x(U^{-1}\grad U)^2\right)
        \right]
    \end{aligned}
\end{equation}
where the two distinct diagonal matrices, $P_\tau$ and $P_x$ presupposes the existence of
three different spin-wave velocities of the problem as obtained from a spin-wave expansion
around the ordered state.



\subsubsection{Multi-parameter RG}

\section{Bond-disordered non-collinear antiferromagnet}

\subsubsection{Random field systems and the replica approach}

\subsubsection{Goldstone bosons and bond disorder}

In our discussion, we closely follow the case of the triangular lattice heisenberg 
antiferromagnet with bond disorder. However, the treatment is in principle easily extensible 
to other spin-models with classicnl non-collinear ordering. We first start by considering 
the effect of an isolated bond impurity. On a non-collinear order this immidiately leads to
local reduction of frustration amongst the spins attached to the bond and therefore, as a 
result, spins reorient. A simple spin-wave expansion around the $120^\circ$ classically
ordered state of the triangular lattice heisenberg antiferromagnet with an isolated bond
impurity leads to the addition of a local perturbing term in the hamiltonian.

\begin{equation}
    \begin{aligned}
        \delta H_i = w c \left(\frac{\delta J\mbf{e}_{\bsym{\delta}}}{J\sqrt{S}a}\right)
        \cdot\grad S_i^\perp
    \end{aligned}
\end{equation}
Where $c = JSa^2$ is the spin-wave velocity, $w$ is a numerical prefactor, 
$\mathbf{e}_{\bsym{\delta}}$ is an unit vector directed across the defective bond and $S^\perp$
is the transverse component of the spin in the body frame of the original spin configuration
and projected onto the ordering plane. The three goldstone modes of the problem emerges
once we expand the term around the three momenta, $Q=(0,0)$ and $Q=(\pm4\pi/3,0)$. It is 
easy to verify that this particular impurity term is most strongly coupled to the in-plane
$Q=(0,0)$ goldstone mode. If a spatial distribution of such impurities are considered we 
arrive at the scenario of bond disorder. With the effect of a single impurity clearly 
understood we can immediately write down a simple free goldstone mode action for the 
disordered case.

\begin{equation}
    \begin{aligned}
        S = \int d\tau\int d^d x\rho\left[\frac{1}{2}\left(\partial_\mu\varphi\right)^2 
        + \mbf{h}\cdot\grad\varphi\right]
    \end{aligned}
\end{equation}
Here $\mbf{h}(\mbf{x})$ are three-component uncorrelated random fields, 
$\overline{\mbf{h}(\mbf{x})}=0$ and $\overline{h_i(\mbf{x})
h_j(\mbf{x}')}=\Delta^2\delta_{ij}\delta^{(d)}(\mbf{x}-\mbf{x}')$. The effect of bond disorder
is therefore to introduce a random field disorder which couples to one of the goldstone 
modes of the symmetry broken phase. Another point to note is that the overall sign of this
additional term is fixed by the chirality symmetry breaking of the ground state. The breakdown
of time reversal symmetry, chirality and the global spin rotation symmetry all conspires
together to generate this random dipolar term. Given, we have a theory of non-interacting
bosons coupled to a random dipolar field, we are still in a position to treat the effect of
disorder exactly. In momentum space the action takes up the following form.

\begin{equation}
    \begin{aligned}
        S =& \rho\int \frac{d\omega}{2\pi}\int \frac{d^dk}{(2\pi)^d}
        \Big[\frac{1}{2}\varphi(-\omega,-\mbf{k})\left(\omega^2+\mbf{k}^2\right)
        \varphi(\omega,\mbf{k})\\
        &+
        2\pi i\delta(\omega)
        \frac{\mbf{h}(-\mbf{k})\cdot\mbf{k}\phi(\omega,\mbf{k})
        -\mbf{h}(\mbf{k})\cdot\mbf{k}\phi(-\omega,-\mbf{k})
        }{2}
        \Big]
    \end{aligned}
\end{equation}
It is easily brought in to a gaussian form by performing the transformation,

\begin{equation}
    \begin{aligned}
        \varphi(\omega,\mbf{k})=\eta(\omega,\mbf{k})-
        2\pi i\delta(\omega)\frac{\mbf{h}(-\mbf{k})\cdot\mbf{k}}{\omega^2+\mbf{k}^2}
    \end{aligned}
\end{equation}
with the action sans an unimportant constant given by, $S =\rho\int \frac{d\omega}{2\pi}
\int \frac{d^dk}{(2\pi)^d}\eta(-\omega,-\mbf{k})\left(\omega^2+\mbf{k}^2\right)
\eta(\omega,\mbf{k})$. The disorder averaged fluctuation for the goldstone modes therefore
reads,
\begin{equation}
    \begin{aligned}
        \overline{\langle\varphi(\tau,\mbf{x})\varphi(\tau',\mbf{x}')\rangle}
        =&\int \frac{d\omega}{2\pi}\int \frac{d^dk}{(2\pi)^d}
        e^{-i\left(\omega(\tau-\tau')+\mbf{k}\cdot(\mbf{x}-\mbf{x}')\right)}
        \langle\eta(-\omega,-\mbf{k})\eta(\omega,\mbf{k})\rangle \\
        &-\Delta^2\int \frac{d^dk}{(2\pi)^d}\frac{1}{\mbf{k}^2}
        e^{-i\left(\mbf{k}\cdot(\mbf{x}-\mbf{x}')\right)}\\
        &=\int \frac{d^Dk}{(2\pi)^D}
        e^{-i\left(k(x-x')\right)}
        \frac{1}{k^2}-\Delta^2\int \frac{d^dk}{(2\pi)^d}\frac{1}{\mbf{k}^2}
        e^{-i\left(\mbf{k}\cdot(\mbf{x}-\mbf{x}')\right)}\\
    \end{aligned}
\end{equation}
with $D=d+1$. Evidently, we see that for non-zero disorder strength $\Delta$, the correlation
funtion has an infrared divergence at $d=2$. This would imply that the goldstone 
modes associated with the non-collinearly ordered triangular lattice antiferromagnet is 
unstable against even very small bond disorder at $d = 2$. In other words, below or at the
critical dimension $d=2$ bond disorder prevents spontaneous symmetry breaking in this
particular antiferromagnet at $T=0$.

Although this argument for bond disorder destroying long range order is very robust, it is 
indeed based on the fluctuation estimates for free goldstone bosons. However, interactions
are important ingradients for an effective theory at low dimensions and they may in 
principle renormalize the disorder coupling of the problem to irrelevancy. Therefore, in 
the following we reconsider the interacting goldstone boson theory for the ground state
of a non-collinear magnet and this time extend it to include the effects of random
exchange couplings.

\subsubsection{Interacting goldstone modes coupled to a random dipolar field}

In order to find the appropriate goldstone mode action in the presence of bond-disorder in a 
non-collinear magnet we have to go back to the microscopic expression for such a term. 
Considering an isolated impurity within one of the elementary plaquettes and between two
particular sublattices we find that the additional contribution from a bond impurity 
$J_{AB} = J + \delta J$ is given by,
\begin{equation}
    \begin{aligned}
        \delta H &= \delta J \ \mbf{S}_A(\mbf{x}_A)\cdot\mbf{S}_B(\mbf{x}_B)
    \end{aligned}
\end{equation}
. Unlike the scenario involving the bulk, the first non-trivial energy contribution from
goldstone mode excitations would now come from the single gradient term along the directed
unit vector $\mbf{u}_{AB}$ along the bond. Like before, we use the convenient orthonormal 
basis transformation and find the impurity contribution to be of the following form.
\begin{equation}
    \begin{aligned}
        \delta H &= \delta JS^2a \ \mbf{n}_AU^{-1}\cdot(\mbf{u}_{AB}\cdot\grad)U\mbf{n}_B\\
        &= -\frac{\delta JS^2a}{\sqrt{12}}\left(
        \mbf{e}_2 U^{-1}\cdot(\mbf{u}_{AB}\cdot\grad)U \mbf{e}_1
        -\mbf{e}_1 U^{-1}\cdot(\mbf{u}_{AB}\cdot\grad)U\mbf{e}_2
        \right)
    \end{aligned}
\end{equation}
where we have exploited the orthonormality of the basis to drop the terms $\mbf{e}_1U^{-1}\grad U
\mbf{e}_1$ and $\mbf{e}_2 U^{-1}\grad U\mbf{e}_2$ which are zero. Similar considerations for bond 
impurity across the other two bond direction reveals that we have the same form of the impurity 
coupling once expressed in terms of the orthonormal basis. Now we introduce a local dipolar field,
$\mbf{h}=\frac{\delta J}{\sqrt{3}Ja}\mbf{u}_{AB}$ and the impurity term assumes the form,
\begin{equation}
    \begin{aligned}
        \delta H &= \mbf{h}\cdot\Tr\left(
        Pt^{3}
        \begin{pmatrix}
            \mbf{e}_1^T\\
            \mbf{e}_2^T\\
            \mbf{e}_3^T
        \end{pmatrix}
        U^{-1}
        \grad
        U
        (\mbf{e}_1,\mbf{e}_2,\mbf{e}_3)\right) \\
        &= \mbf{h}\cdot\Tr\left(Pt^3 U^{-1}\grad U\right)
    \end{aligned}
\end{equation}
where the same gauge choice from the bulk action was adopted and it is to be noticed that the 
$SO(3)$ generator corresponding to in plane rotation 
$t^3=\begin{pmatrix}0&1&0\\-1&0&0\\0&0&0\end{pmatrix}$ makes an appearance. The latter has to do 
with the fact that in the triangular lattice the local dipolar field arising due to bond disorder 
couples directly with the goldstone mode associated with the in plane rotation. Now, it is 
straightforward to generalize to the case of spatially distributed bond impurities i.e. disordered
bonds. For a generic scenario involving the emergent random dipolar field coupled with multiple 
goldstone modes the hamiltonian shall take up the following form.
\begin{equation}
    \begin{aligned}
        H = \int d^dx\left[
            \Tr\left(P(U^{-1}\grad U)^2\right)
            +\mbf{h}\cdot\Tr\left(P(c_a t^a)(U^{-1}\grad U)\right)
        \right]
    \end{aligned}
\end{equation}
As before the interacting nature of the goldstone modes are apparent once we perform an expansion
of the form $U^{-1}\partial_\mu U$ in terms of the $G/H$ generators with 
$U=e^{\bsym{\xi}\cdot\mbf{t}}$.
\begin{equation}
    \begin{aligned}
        U^{-1}\partial_\mu U = t^a\omega^\mu_a(\bsym{\xi})
        = t^a\left[
            \partial_\mu\xi_a
            -\frac{1}{2}f_{abc}\xi_b\partial_\mu\xi_c
            +\frac{1}{6}f_{abc}f_{cde}\xi_b\xi_d\partial_\mu\xi_e
            +O(\xi^4)
            \right]
    \end{aligned}
\end{equation}
Since, we shall not be concerned with the critical point of the ordered phase we can ignore the 
differences between various spin-wave velocities and obtain a simpler goldstone mode action.
\begin{equation}
    \begin{aligned}
        S &= \frac{1}{2}\int d\tau\int d^dx\left[
            \rho_a\omega^\mu_a(\tau,\mbf{x})\omega^\mu_a(\tau,\mbf{x})
            +c_a\rho_ah_i(\mbf{x})\omega^i_a(\tau,\mbf{x})
        \right]
    \end{aligned}
\end{equation}
where we used the normalization $\Tr(Pt^at^b)=(\rho_a/2)\delta^{ab}$ and $f_{abc}$ are the 
structure constants of the broken symmetry $G/H$,
\begin{equation}
    \begin{aligned}
        \left[t^a,t^b\right]=f_{abc}t^c
    \end{aligned}
\end{equation}
. In the case of triangular lattice we have $ G/H = SO(3)$ so $f_{abc}=\epsilon_{abc}$, i.e. 
the levi-civita tensor.

With random bond disorder we may consider a gaussian distribution for the inhomogeneous dipolar 
fields,
\begin{equation}
    \begin{aligned}
        P\left[\mbf{h}\right]=\int\mathcal{D}\mbf{h} \ 
        e^{-\frac{1}{2\Delta^2}\int d^dx\mbf{h}(\mbf{x})^2}
    \end{aligned}
\end{equation}
and this leads to the homogeneous replica action of our interest.
\begin{equation}
    \begin{aligned}
        S =& \frac{1}{2}\Big[
            \rho_a
            \sum_{r}\int d\tau\int d^dx\
            \omega^{r,\mu}_a(\tau,\mbf{x})\omega^{r,\mu}_a(\tau,\mbf{x})\\
            &-
            c^2_a\rho^2_a\Delta^2
            \sum_{r,s}
            \int d\tau\int d\tau'\int d^dx\
            \omega^{r,i}_a(\tau,\mbf{x})
            \omega^{s,i}_a(\tau',\mbf{x})
            \Big]
    \end{aligned}
\end{equation}

There are two different types of couplings of the problem, $\rho_a$ and $\Delta$ and their mass 
dimensions are the following.
\begin{equation}
    \begin{aligned}
        \left[\rho_a\right] &= d - 1\\
        \left[\Delta\right] &= (2-d)/2
    \end{aligned}
\end{equation}
This immediately tells us that disorder is marginal in $d=2$ and relevant at $d<2$. This
result is consistent with the dimensionality of the infrared catastrophe related to the
goldstone mode fluctuation. However, at the same time we face the problem that our
homogeneous couplings and the disorder coupling do not have a well defined perturbative 
expansion at the same spatial dimension for $T=0$. Around $d=1$ the homogeneous couplings
$\rho_a$ are marginal but the disorder coupling $\Delta$ is strongly relevant and flows to
infinity. Unfortunately this does not tell us anything about what happens close to spatial 
dimension $d=2$ when the disorder is marginal. Therefore, in the following we consider the 
alternate scenario with an expansion at $d=2-\epsilon$. The homogeneous coupling is relevant 
now but we shall only be concerned with the independent flow of the disorder coupling.

\subsubsection{Renormalization of the disorder coupling $\Delta$ at $d=2-\epsilon$}

\begin{thebibliography}{99}

    \bibitem{imry-ma-75}
        Imry and Ma, 75

    \bibitem{mermin-wagner-66}
        Mermin and Wagner, 66

    \bibitem{aharony-78}
        Aharony, 78

    \bibitem{azaria-92}
        Azaria, Delamotte and Mouhanna, 92

    \bibitem{azaria-93}
        Azaria, Delamotte, Delduc and Jolicoeur, 93

    \bibitem{dombre-88}
        Dombre and Read, 88

\end{thebibliography}

\end{document}
